\chapter{Ruang Hilbert}\label{ch:b3:hilbert}

\omterm{def:b3:hilbert:inner}{Ruang Hilbert} adalah ruang Banach yang normanya datang dari
sebuah \omterm{def:b3:hilbert:inner}{hasil kali dalam} --- dan satu struktur tambahan itu memulihkan, dalam
dimensi tak berhingga, hampir seluruh geometri Euclid:
proyeksi ortogonalnya ada, setiap fungsional \omterm{def:b3:topology:continuity}{kontinu} berupa
\omterm{def:b3:hilbert:inner}{hasil kali dalam} terhadap sebuah vektor tetap (menurut Riesz), dan
basis ortonormalnya menguraikan setiap vektor menjadi deret konvergen dengan
pembukuan Pythagoras (menurut \omterm{thm:b3:hilbert:parseval}{Parseval}). Puncak bab ini adalah sebuah
utang yang dilunasi: bahwa sistem trigonometri merupakan basis ortonormal
$L^2$, sehingga identitas \omterm{thm:b3:hilbert:parseval}{Parseval} berlaku bagi \emph{setiap}
fungsi yang \omterm{def:b3:lebesgue:l1}{terintegralkan} kuadrat --- yakni pernyataan yang pada Tahun ke-2 hanya
dapat dibuktikan untuk fungsi $\mathcal C^1$ sepenggal-sepenggal. Kita mengakhirinya dengan
Lax--Milgram, yakni lema kuda beban pendekatan variasi
terhadap persamaan diferensial.

Sepanjang bab ini, $H$ adalah ruang vektor atas $K = \R$ atau $\C$.

\section{Hasil kali dalam; teorema proyeksi}

\begin{definition}\label{def:b3:hilbert:inner}
Sebuah \emph{hasil kali dalam}\index{hasil kali dalam} adalah pemetaan $\langle
\cdot,\cdot\rangle \colon H\times H \to K$, yang linear pada
variabel keduanya, dengan $\langle y, x\rangle =
\overline{\langle x, y\rangle}$ dan $\langle x, x\rangle > 0$
untuk $x \neq 0$. Ia menginduksi norma $\norm x = \langle x,
x\rangle^{1/2}$, \emph{ketaksamaan Cauchy--Schwarz}
$\abs{\langle x, y\rangle} \leq \norm x\norm y$ (dengan bukti Tahun
ke-2 --- lewat diskriminannya --- yang tak berubah), dan
\emph{hukum jajaran genjang}
\[
\norm{x + y}^2 + \norm{x - y}^2 = 2\norm x^2 + 2\norm y^2 .
\]
Sebuah \emph{ruang Hilbert}\index{ruang Hilbert} adalah ruang
hasil kali dalam yang \omterm{def:b3:complete:complete}{lengkap} terhadap norma ini. Contohnya: $\ell^2$
(\cref{pb:b3:banach:1}) dan, yang paling mendasar, $L^2(\mu)$
dengan $\langle f, g\rangle = \int\bar fg\,\dd\mu$ ---
yang \omterm{def:b3:complete:complete}{lengkap} menurut Riesz--Fischer (\cref{thm:b3:lp:complete}); sedangkan
hasil kali dalamnya berhingga menurut Cauchy--Schwarz ($=$ Hölder pada
$p = q = 2$).
\end{definition}

\begin{theorem}[Proyeksi ke sebuah himpunan cembung tertutup]
\label{thm:b3:hilbert:projection}
Misalkan $C \neq \varnothing$ sebuah himpunan bagian \emph{cembung} tertutup dari
\omterm{def:b3:hilbert:inner}{ruang Hilbert} $H$ dan $x \in H$. Maka ada tepat satu $p_C(x)
\in C$ dengan\index{teorema proyeksi}
\[
\norm{x - p_C(x)} = d(x, C),
\]
yang dicirikan oleh: $\operatorname{Re}\langle x - p_C(x),\ c -
p_C(x)\rangle \leq 0$ untuk setiap $c \in C$. Dan pemetaan $p_C$ bersifat
Lipschitz-$1$.
\end{theorem}

\begin{proof}
Misalkan $d = d(x, C)$ dan $(c_n) \subseteq C$ dengan $\norm{x - c_n}
\to d$. Jajaran genjang pada $x - c_n$ dan $x - c_m$:
\[
\norm{c_n - c_m}^2 = 2\norm{x - c_n}^2 + 2\norm{x - c_m}^2 -
4\,\bigl\|x - \tfrac{c_n + c_m}2\bigr\|^2
\leq 2\norm{x{-}c_n}^2 + 2\norm{x{-}c_m}^2 - 4d^2
\]
(sebab kecembungannya menaruh titik tengahnya di $C$): jadi ruas kanannya menuju
$0$, sehingga $(c_n)$ Cauchy, dan limitnya $p \in C$ (yang tertutup)
mencapai $d$. Ketunggalannya: dua pemberi minimum memberikan, lewat identitas
yang sama, $\norm{p - p'}^2 \leq 2d^2 + 2d^2 - 4d^2 = 0$.

Pencirinya: untuk $c \in C$ dan $t \in \intoc01$, vektor
$p + t(c - p) \in C$, sehingga
\[
d^2 \leq \norm{x - p - t(c-p)}^2
= d^2 - 2t\operatorname{Re}\langle x - p, c - p\rangle +
t^2\norm{c-p}^2 ;
\]
lalu bagilah dengan $t \to 0^+$: $\operatorname{Re}\langle x - p, c -
p\rangle \leq 0$. Sebaliknya ketaksamaan ini memberikan $\norm{x -
c}^2 = \norm{x - p}^2 - 2\operatorname{Re}\langle x - p, c -
p\rangle + \norm{p - c}^2 \geq \norm{x-p}^2$. Untuk kelipschitzannya: bagi
$x, y$ yang berproyeksi $p, q$, jumlahkanlah kedua
ketaksamaan variasinya (dengan $c = q$, berturut-turut $c = p$):
$\operatorname{Re}\langle x - y - (p - q), p - q\rangle \geq
0$, sehingga $\norm{p - q}^2 \leq \operatorname{Re}\langle x - y, p -
q\rangle \leq \norm{x - y}\norm{p - q}$.
\end{proof}

\begin{theorem}[Penguraian ortogonal]
\label{thm:b3:hilbert:decomposition}
Misalkan $F$ sebuah \emph{subruang tertutup} dari $H$. Maka $p_F$ bersifat
linear, $x - p_F(x) \perp F$ untuk setiap $x$, dan\index{komplemen
ortogonal}
\[
H = F \oplus F^\perp,
\qquad F^\perp = \{y : \langle y, f\rangle = 0\ \forall f\in
F\},
\qquad (F^\perp)^\perp = F .
\]
Untuk subruang yang umum, $(F^\perp)^\perp = \bar F$; khususnya
$F$ padat jika dan hanya jika $F^\perp = \{0\}$.
\end{theorem}

\begin{proof}
Untuk sebuah subruang, pencirian variasinya dengan $c =
p_F(x) \pm f$ (dengan $f \in F$, kedua tandanya, dan $\iu f$ pada
kasus kompleksnya) memaksa $\langle x - p_F(x), f\rangle = 0$: jadi
sisanya ortogonal terhadap $F$. Penguraiannya $x = p_F(x) + (x
- p_F(x))$ dengan $F \cap F^\perp = \{0\}$ (sebab $\langle y, y
\rangle = 0$); sedangkan kelinearan $p_F$ menyusul dari ketunggalan
penguraian semacam itu (sebab kedua ruasnya linear di dalamnya). Selalu $(F^\perp)
^\perp \supseteq F$; sebaliknya jika $x \perp F^\perp$,
tulislah $x = f + g$: maka $g = x - f \in F^\perp$ dan $\langle g,
g\rangle = \langle x, g\rangle - \langle f, g\rangle = 0$: jadi $x
= f \in F$. Untuk subruang $F$ yang umum: $F^\perp = \bar
F^{\,\perp}$ (lewat \omterm{def:b3:topology:continuity}{kekontinuan} \omterm{def:b3:hilbert:inner}{hasil kali dalamnya}), sehingga
$(F^\perp)^\perp = \bar F$ menurut kasus tertutupnya; dan \omterm{ex:b3:lebesgue:gamma}{kepadatannya} jika dan hanya jika
$\bar F = H$ jika dan hanya jika $F^\perp = 0$.
\end{proof}

\begin{theorem}[Penyajian Riesz]\label{thm:b3:hilbert:riesz}
Untuk setiap fungsional linear \omterm{def:b3:topology:continuity}{kontinu} $\varphi \in H'$ ada
tepat satu $a \in H$ dengan\index{teorema penyajian Riesz}
\[
\varphi(x) = \langle a, x\rangle \quad (x \in H),
\qquad \norm\varphi_{H'} = \norm a .
\]
\end{theorem}

\begin{proof}
Jika $\varphi = 0$: maka $a = 0$. Selain itu $F = \ker\varphi$ merupakan
subruang sejati yang tertutup; lalu pilihlah $u \in F^\perp$ dengan $\norm u = 1$
(menurut \cref{thm:b3:hilbert:decomposition}: sebab $F^\perp \neq 0$ karena
$F \neq H$). Untuk sembarang $x$, vektor $\varphi(x)u -
\varphi(u)x \in \ker\varphi$, sehingga $\perp u$:
\[
0 = \langle u, \varphi(x)u - \varphi(u)x\rangle
= \varphi(x) - \varphi(u)\langle u, x\rangle :
\qquad \varphi(x) = \langle
\overline{\varphi(u)}\,u,\ x\rangle .
\]
Jadi $a = \overline{\varphi(u)}u$ berhasil. Ketunggalannya: $\langle a
- a', x\rangle = 0$ untuk setiap $x$, lalu ujilah $x = a - a'$. Normanya:
$\abs{\varphi(x)} \leq \norm a\norm x$ (lewat Cauchy--Schwarz) dengan
kesamaannya di $x = a$.
\end{proof}

\begin{example}[Sebuah proyeksi, yang dihitung sampai tuntas]
\label{ex:b3:hilbert:projectionexample}
Di $H = L^2(\intcc01)$, apakah penghampiran terbaik atas
$f(x) = x^2$ oleh fungsi afin? Subruang $F =
\operatorname{Vect}(1, x)$ bersifat tertutup (sebab berdimensi berhingga),
dan $p_F(f) = a + bx$ dicirikan oleh keortogonalan
sisanya terhadap $1$ dan terhadap $x$:
\[
\int_0^1(x^2 - a - bx)\,\dd x = 0,
\qquad
\int_0^1x\,(x^2 - a - bx)\,\dd x = 0,
\]
yakni $\frac13 = a + \frac b2$ dan $\frac14 = \frac a2 +
\frac b3$: sehingga $a = -\frac16$ dan $b = 1$. Jadi $p_F(x^2) = x -
\frac16$, dan galatnya adalah
\[
d(f, F)^2 = \int_0^1\Bigl(x^2 - x + \frac16\Bigr)^2\dd x =
\frac1{180},
\qquad d(f, F) = \frac1{6\sqrt5} .
\]
Dua catatan yang layak diendapkan. Pertama, perhitungannya
tak lain sistem linear $2\times2$ --- yakni \emph{persamaan
normalnya}; dan untuk basis monomialnya matriksnya
$\bigl(\frac1{i+j+1}\bigr)$ adalah matriks Hilbert yang
terkenal buruk kondisinya, sedangkan mengortogonalkannya lebih dahulu
(lewat polinomial Legendre, \cref{pb:b3:hilbert:1}) merupakan obatnya.
Kedua, penghampiran \emph{seragam} terbaik atas $x^2$ oleh
fungsi afin berbeda ($x - \frac18$, lewat
osilasi setara): sebab setiap norma mempunyai geometrinya sendiri, dan hanya
yang Hilbertlah yang menjawab dengan sebuah sistem linear.
\end{example}

\section{Basis ortonormal}

\begin{definition}\label{def:b3:hilbert:onb}
Sebuah keluarga $(e_i)_{i\in I}$ disebut \emph{ortonormal} jika $\langle
e_i, e_j\rangle = \delta_{ij}$, dan sebuah \emph{basis
Hilbert}\index{basis Hilbert} (yakni basis ortonormal) jika lebih lanjut
kombinasi linear berhingganya padat di $H$ (yakni keluarganya
bersifat \emph{total}). Kita menangani kasus terbilangnya $I = \N$, yang lewat
Gram--Schmidt mencakup setiap $H$ yang \emph{terpisahkan}
(\cref{prop:b3:hilbert:gramschmidt}).
\end{definition}

\begin{theorem}[Bessel, Parseval]\label{thm:b3:hilbert:parseval}
Misalkan $(e_n)_{n\in\N}$ ortonormal di $H$, dan $c_n(x) =
\langle e_n, x\rangle$.
\begin{enumerate}
  \item (Bessel) $\sum_n\abs{c_n(x)}^2 \leq \norm x^2$, dan
        deret $\sum_nc_n(x)e_n$ konvergen di $H$, dengan
        jumlah $p_F(x)$ dan $F = \overline{\operatorname{Vect}}(e_n)$.
  \item Pernyataan berikut setara: (i) $(e_n)$ merupakan \omterm{def:b3:hilbert:onb}{basis
        Hilbert}; (ii) $x = \sum_nc_n(x)e_n$ untuk setiap $x$; (iii)
        \emph{Parseval}\index{identitas Parseval}: $\norm x^2
        = \sum_n\abs{c_n(x)}^2$ untuk setiap $x$; (iv) satu-satunya
        vektor yang ortogonal terhadap semua $e_n$ adalah $0$.
  \item Jika $(e_n)$ \omterm{def:b3:hilbert:onb}{basis Hilbert}, maka $x \mapsto (c_n(x))_n$
        merupakan isomorfisma isometrik $H \to \ell^2$
        (sehingga \emph{setiap} \omterm{def:b3:hilbert:inner}{ruang Hilbert} terpisahkan berdimensi
        tak berhingga ``adalah'' $\ell^2$), dan $\langle x, y\rangle =
        \sum_n\overline{c_n(x)}c_n(y)$.
\end{enumerate}
\end{theorem}

\begin{proof}
(1) Untuk $N$ yang berhingga: $x - \sum_{n\leq N}c_ne_n \perp e_k$ (bila $k
\leq N$), sehingga Pythagoras memberikan $\norm x^2 = \sum_{n\leq
N}\abs{c_n}^2 + \norm{x - \sum_{n\leq N}c_ne_n}^2$: yakni Bessel.
Lalu jumlah parsialnya $S_N = \sum_{n\leq N}c_ne_n$ bersifat Cauchy:
sebab $\norm{S_N - S_M}^2 = \sum_{M<n\leq N}\abs{c_n}^2$, yakni ekor sebuah
deret konvergen; limitnya terletak di $F$, dan $x - \lim S_N
\perp$ setiap $e_k$ (lewat \omterm{def:b3:topology:continuity}{kekontinuannya}), jadi $\perp F$: sehingga menurut ketunggalan
penguraian ortogonalnya, $\lim S_N = p_F(x)$.

(2) (i)$\Rightarrow$(ii): sebab $F = H$, sehingga $p_F = \mathrm{id}$.
(ii)$\Rightarrow$(iii): lewat Pythagoras pada limitnya (sebab $\norm{S_N}^2
= \sum_{n \leq N}\abs{c_n}^2 \to \norm x^2$).
(iii)$\Rightarrow$(iv): sebab $x \perp$ semua $e_n$ memberikan $\norm x^2 =
0$. (iv)$\Rightarrow$(i): sebab $F^\perp = \{0\}$
(karena keortogonalan terhadap semua $e_n$ adalah keortogonalan terhadap $F$), sehingga $F$
padat menurut \cref{thm:b3:hilbert:decomposition}; padahal $F$, sebagai
sebuah \omterm{def:b3:topology:interior}{penutup}, sudah tertutup: jadi $F = H$.

(3) Pemetaannya linear, isometrik menurut (iii) (sehingga injektif),
dan surjektif: sebab diberikan $(c_n) \in \ell^2$, deret $\sum
c_ne_n$ konvergen (yang Cauchy seperti pada (1)) ke sebuah prapetanya. Sedangkan rumus
\omterm{def:b3:hilbert:inner}{hasil kali dalamnya} berupa polarisasi dari (iii), atau perhitungan limit
yang langsung.
\end{proof}

\begin{proposition}[Gram--Schmidt]\label{prop:b3:hilbert:gramschmidt}
Misalkan $(x_n)$ sebuah barisan yang bebas linear. Dengan menetapkan
secara induktif $\tilde e_n = x_n - \sum_{k<n}\langle e_k,
x_n\rangle e_k$ dan $e_n = \tilde e_n/\norm{\tilde e_n}$
dihasilkanlah $(e_n)$ yang ortonormal dengan rentang berhingga yang sama:
$\operatorname{Vect}(e_1, \dots, e_n) = \operatorname{Vect}
(x_1, \dots, x_n)$. Karena itu setiap \omterm{def:b3:hilbert:inner}{ruang Hilbert} terpisahkan
(yakni yang mempunyai himpunan bagian padat terbilang) mempunyai \omterm{def:b3:hilbert:onb}{basis
Hilbert}.\index{pengortonormalan Gram--Schmidt}
\end{proposition}

\begin{proof}
Lewat induksi: $\tilde e_n \perp e_k$ (bila $k < n$) menurut konstruksinya,
dan $\tilde e_n \neq 0$ menurut kebebasannya; sedangkan rentangnya bersesuaian pada
setiap tahapnya (lewat penggantian basis segitiga). Untuk $H$ yang terpisahkan:
dari sebuah barisan padat petiklah subkeluarga yang bebas
linear dengan rentang yang padat (buanglah setiap vektor yang berada di rentang
pendahulunya, sebab rentangnya tak berubah), lalu ortonormalkan:
sehingga hasilnya total.
\end{proof}

\begin{theorem}[Sistem trigonometri; Parseval akhirnya]
\label{thm:b3:hilbert:fourier}
Di $L^2(\intcc{-\pi}\pi)$ dengan $\langle f, g\rangle =
\frac1{2\pi}\int_{-\pi}^\pi \bar fg$, keluarga $e_n(t) =
\eu^{\iu nt}$ dengan $n \in \Z$ merupakan \omterm{def:b3:hilbert:onb}{basis Hilbert}. Karena itu,
untuk \emph{setiap} $f \in L^2$ --- khususnya setiap fungsi
$2\pi$-periodik $f$ yang \omterm{def:b3:topology:continuity}{kontinu} sepenggal-sepenggal --- dengan $c_n(f) =
\frac1{2\pi}\int_{-\pi}^{\pi}f(t)\eu^{-\iu nt}\dd t$:
\[
f = \sum_{n\in\Z}c_n(f)\,\eu^{\iu nt} \ \ \text{di } L^2,
\qquad
\frac1{2\pi}\int_{-\pi}^{\pi}\abs f^2 =
\sum_{n\in\Z}\abs{c_n(f)}^2 .
\]
Ini membuktikan, dalam keumuman penuh, identitas \omterm{thm:b3:hilbert:parseval}{Parseval} yang
diterima begitu saja pada Tahun ke-2.\index{identitas Parseval (umum)}
\end{theorem}

\begin{proof}
Keortonormalannya berupa perhitungan langsung (Tahun ke-2). Untuk ketotalannya: misalkan
$f \in L^2$ bersifat $\perp$ semua $e_n$, yakni semua koefisien
Fouriernya lenyap. Fungsi \omterm{def:b3:topology:continuity}{kontinu} berperiode $2\pi$ bersifat
padat di $L^2(\intcc{-\pi}\pi)$: sebab sungguh $\mathcal
C_c(\intoo{-\pi}\pi)$ padat
(\cref{thm:b3:lp:density}(2)) dan fungsi semacam itu diperluas
secara periodik dan \omterm{def:b3:topology:continuity}{kontinu}. Sedangkan polinomial trigonometri bersifat
padat terhadap $\norm\cdot_\infty$ di antara fungsi periodik \omterm{def:b3:topology:continuity}{kontinu}
(menurut Stone--Weierstrass, \cref{cor:b3:complete:weierstrass}(c)),
dan $\norm\cdot_2 \leq \norm\cdot_\infty$: sehingga polinomial
trigonometrinya padat di $L^2$. Padahal $f \perp$ setiap
polinomial trigonometri, jadi $f \perp$ sebuah subruang padat: sehingga $f
\in (\text{padat})^\perp = \{0\}$
(\cref{thm:b3:hilbert:decomposition}). Lalu kriteria (iv)
\cref{thm:b3:hilbert:parseval} menuntaskannya; sedangkan (ii) dan (iii)
terurai menjadi tampilannya (dengan mengindeks ulang $\Z$ yang terbilang; dan
deret dua ujungnya konvergen tanpa syarat --- sebab semua jumlah
parsialnya atas sembarang keluarga yang menghabiskan indeksnya konvergen, lewat argumen ekor $\ell^2$-nya).
\end{proof}

\begin{theorem}[Lax--Milgram]\label{thm:b3:hilbert:laxmilgram}
Misalkan $H$ \omterm{def:b3:hilbert:inner}{ruang Hilbert} real dan $a \colon H\times H \to
\R$ bilinear, \emph{\omterm{def:b3:topology:continuity}{kontinu}} (yakni $\abs{a(u,v)} \leq M\norm
u\norm v$) serta \emph{koersif} (yakni $a(u, u) \geq \alpha\norm u^2$
dengan $\alpha > 0$). Maka untuk setiap $\varphi \in H'$ ada
tepat satu $u \in H$ dengan\index{teorema Lax--Milgram}
\[
a(u, v) = \varphi(v) \qquad \text{untuk setiap } v \in H .
\]
\end{theorem}

\begin{proof}
Untuk $u$ yang tetap, $v \mapsto a(u, v)$ bersifat linear \omterm{def:b3:topology:continuity}{kontinu}: sehingga Riesz
memberikan tepat satu $Au \in H$ dengan $a(u,v) = \langle Au,
v\rangle$; dan $A$ linear dengan $\norm{Au} \leq M\norm u$
(lewat ketunggalan wakilnya, lalu batasnya). Untuk kekoersifannya:
$\alpha\norm u^2 \leq a(u,u) = \langle Au, u\rangle \leq
\norm{Au}\norm u$, sehingga $\norm{Au} \geq \alpha\norm u$: jadi $A$
injektif dengan jangkauan yang tertutup (sebab barisan peta Cauchy $Au_n$
memaksa $u_n$ Cauchy). Jangkauannya padat: sebab $w \perp
\operatorname{im}A$ memberikan $0 = \langle Aw, w\rangle \geq
\alpha\norm w^2$. Tertutup dan padat: jadi $A$ bijektif. Diberikan
$\varphi$, misalkan $f$ menyajikannya (lewat Riesz) dan $u = A^{-1}f$:
maka $a(u, v) = \langle f, v\rangle = \varphi(v)$, secara tunggal
(sebab $a(u - u', \cdot) = 0$ dan kekoersifannya).
\end{proof}

\begin{remark}\label{rem:b3:hilbert:variational}
Bila $a$ setangkup, solusi Lax--Milgram merupakan satu-satunya
pemberi minimum \emph{energi} $J(v) = \frac12a(v,v) -
\varphi(v)$ (\cref{exo:b3:hilbert:9}): jadi keberadaan solusi
soal variasi dalam satu sapuan. Diterapkan pada ruang fungsi
yang sesuai (yakni ruang Sobolev pada kuliah berikutnya), ini
menyelesaikan soal nilai batas bagi persamaan diferensial ---
yakni pintu masuk modern ke persamaan diferensial parsial.
\end{remark}

\section{Latihan}

\begin{exercise}[$\star$]\label{exo:b3:hilbert:1}
(a) Buktikan identitas polarisasinya (yang real: $4\langle x,
y\rangle = \norm{x+y}^2 - \norm{x-y}^2$; dan yang kompleks: versi
empat sukunya).
(b) Tunjukkan bahwa $\norm\cdot_1$ pada $L^1(\intcc01)$ dan
$\norm\cdot_\infty$ pada $\mathcal C(\intcc01)$ melanggar
hukum jajaran genjang: sehingga norma ini tak datang dari \omterm{def:b3:hilbert:inner}{hasil kali dalam} mana pun.
\end{exercise}

\begin{exercise}[$\star$]\label{exo:b3:hilbert:2}
Di $H = L^2(\intcc01)$ (yang real):
(a) hitunglah proyeksi $f$ ke subruang berisi
fungsi konstan, lalu tafsirkan;
(b) hitunglah proyeksinya ke $\{g : g = 0 \text{ hampir di mana-mana pada }
\intcc0{1/2}\}$;
(c) hitunglah $d\bigl(x \mapsto x,\ \operatorname{Vect}(\mathbf
1)\bigr)$.
\end{exercise}

\begin{exercise}[$\star\star$]\label{exo:b3:hilbert:3}
(a) Tunjukkan bahwa untuk sebuah subruang $F$: $F$ padat $\iff$ $F^\perp =
\{0\}$, lalu berikan contoh di $\ell^2$ sebuah subruang padat
yang \emph{sejati} (sehingga $F^\perp = 0$ tanpa $F = H$: yakni
teorema penguraiannya sungguh memerlukan $F$ yang tertutup).
(b) Tunjukkan bahwa jika $x_n \to x$ dan $y_n \to y$ dalam norma, maka
$\langle x_n, y_n\rangle \to \langle x, y\rangle$, lalu temukan
dua tempat yang bab ini memakai \omterm{def:b3:topology:continuity}{kekontinuan} itu.
\end{exercise}

\begin{exercise}[$\star\star$]\label{exo:b3:hilbert:4}
Terapkan Gram--Schmidt pada $1, x, x^2$ di $L^2(\intcc{-1}1)$
(dengan \omterm{def:b3:measure:lebesgueouter}{ukuran Lebesgue}): perolehlah ketiga
\emph{polinomial Legendre} ternormalkan yang pertama, lalu periksalah bahwa keduanya cocok dengan
$\sqrt{n + \frac12}\,P_n$ untuk polinomial Rodrigues $P_n$
pada \cref{pb:b3:hilbert:1}.
\end{exercise}

\begin{exercise}[$\star\star$]\label{exo:b3:hilbert:5}
Terapkan \omterm{thm:b3:hilbert:parseval}{Parseval} (\cref{thm:b3:hilbert:fourier}) pada $f(t) = t$
dan $f(t) = t^2$ di $\intcc{-\pi}\pi$ --- kini secara sah bagi
keduanya (yang \omterm{def:b3:topology:continuity}{kontinu}, tetapi sebelumnya identitasnya menuntut
kehati-hatian $\mathcal C^1$ sepenggal-sepenggal pada ketakkontinuan di sambungannya):
pulihkanlah
\[
\sum_{n\geq1}\frac1{n^2} = \frac{\pi^2}6,
\qquad
\sum_{n\geq1}\frac1{n^4} = \frac{\pi^4}{90} .
\]
\end{exercise}

\begin{exercise}[$\star\star$]\label{exo:b3:hilbert:6}
(a) Carilah $a \in L^2(\intcc01)$ dengan $\int_0^{1/2}f =
\langle a, f\rangle$ untuk setiap $f$; lalu hitunglah $\norm\varphi$ untuk
fungsional ini.
(b) Tunjukkan bahwa penilaian $f \mapsto f(\frac12)$, yang terdefinisi
pada subruang $\mathcal C(\intcc01) \subseteq
L^2(\intcc01)$, \emph{tak} \omterm{def:b3:topology:continuity}{kontinu} terhadap
$\norm\cdot_2$: sehingga tak ada wakil Riesznya (sebab penilaian bukan
gagasan $L^2$).
\end{exercise}

\begin{exercise}[$\star\star\star$]\label{exo:b3:hilbert:7}
Misalkan $H$ terpisahkan dengan \omterm{def:b3:hilbert:onb}{basis Hilbert} $(e_n)$, dan $(x_k)$ sebuah
barisan yang terbatas.
(a) Tunjukkan bahwa suatu subbarisannya konvergen \emph{secara lemah}: yakni
ada $x$ dengan $\langle y, x_{k_j}\rangle \to \langle y,
x\rangle$ untuk setiap $y \in H$. \emph{(Pakai pemetikan diagonal pada
koefisiennya $\langle e_n, x_k\rangle$; lalu rakitlah $x$ lewat
Bessel dan keterbatasan seragam normanya.)}
(b) Tunjukkan bahwa $e_n \rightharpoonup 0$ padahal $\norm{e_n} = 1$: sehingga limit
lemah dapat kehilangan norma. Lalu tunjukkan $\norm x \leq
\liminf\norm{x_{k_j}}$ pada (a).
\end{exercise}

\begin{exercise}[$\star\star$]\label{exo:b3:hilbert:8}
(Adjoin) Untuk $T \in \mathcal L(H)$, tunjukkan bahwa ada tepat satu
$T^* \in \mathcal L(H)$ dengan $\langle Tx, y\rangle = \langle
x, T^*y\rangle$ (lewat Riesz), dan $\vertiii{T^*} = \vertiii T$.
Hitunglah adjoin pergeseran $S$ pada $\ell^2$, lalu buktikan
$\ker T^* = (\operatorname{im}T)^\perp$ --- lalu turunkan
$\overline{\operatorname{im}T} = (\ker T^*)^\perp$.
\end{exercise}

\begin{exercise}[$\star\star$]\label{exo:b3:hilbert:9}
Misalkan $a$ seperti pada Lax--Milgram dan lebih lanjut \emph{setangkup}.
Tunjukkan bahwa $u$ menyelesaikan $a(u, \cdot) = \varphi$ jika dan hanya jika $u$ meminimumkan
$J(v) = \frac12a(v, v) - \varphi(v)$, dan bahwa minimumnya
tercapai tepat di satu titik. \emph{(Lengkapkan kuadratnya:
$J(u + w) - J(u) = \frac12a(w,w) \geq
\frac\alpha2\norm w^2$.)} Terapannya: turunkan kembali
teorema proyeksinya bagi subruang tertutup dari Lax--Milgram.
\end{exercise}

\begin{exercise}[$\star\star\star$]\label{exo:b3:hilbert:10}
(Sistem Haar) Pada $\intcc01$, misalkan $h_{0} = \mathbf 1$, dan
untuk $n = 2^j + k$ (dengan $j \geq 0$ dan $0 \leq k < 2^j$):
\[
h_n = 2^{j/2}\Bigl(\mathbf 1_{[k2^{-j},\,(k +
\frac12)2^{-j})} - \mathbf 1_{[(k+\frac12)2^{-j},\,(k+1)2^{-j})}
\Bigr).
\]
Tunjukkan bahwa $(h_n)_{n\geq0}$ ortonormal di $L^2(\intcc01)$,
dan total. \emph{(Keortogonalannya: sebab pendukungnya saling lepas atau bersarang;
sedangkan ketotalannya: sebab rentang berhingganya memuat semua fungsi tangga diadik,
yang padat --- lewat \cref{thm:b3:lp:density}(1) dan penghampiran
diadik atas selangnya.)} Sistem Haar merupakan nenek moyang
gelombang kecil.
\end{exercise}

\begin{exercise}[$\star\star$]\label{exo:b3:hilbert:11}
(Proyeksi ortogonal, yang tercirikan) Misalkan $H$ sebuah \omterm{def:b3:hilbert:inner}{ruang
Hilbert} dan $P \in \mathcal L(H)$ dengan $P^2 = P$ dan $P \neq 0$.
Tunjukkan kesetaraan:
(i) $P$ merupakan proyeksi ortogonal ke
$\operatorname{im}P$;
(ii) $P = P^*$ (\cref{exo:b3:hilbert:8});
(iii) $\vertiii P = 1$.
\emph{(Untuk (iii) $\Rightarrow$ (i): jika suatu $x \in
(\ker P)^\perp$ mempunyai $Px \neq x$, tinjaulah $x + t(Px - x)$
--- atau secara langsung: untuk $u \in \operatorname{im}P$ dan $v \in
\ker P$, uraikan $\norm{P(u + tv)}^2 \leq \norm{u + tv}^2$
untuk setiap $t \in \R$ lalu simpulkan $\langle u, v\rangle = 0$.)}
Tampilkan sebuah proyeksi tak ortogonal pada $\R^2$ lalu hitunglah
normanya.
\end{exercise}

\begin{exercise}[$\star\star\star$]\label{exo:b3:hilbert:12}
(Teorema ergodik von Neumann) Misalkan $U \in \mathcal L(H)$
\emph{uniter} (yakni $U^*U = UU^* = I$), $F = \ker(U - I)$
ruang tetapnya, $P$ proyeksi ortogonal ke $F$, dan
$A_n = \frac1n\sum_{k=0}^{n-1}U^k$.
(a) Tunjukkan bahwa $\ker(U - I) = \ker(U^* - I)$ \emph{(dari
$\norm{Ux - x}^2 = 2\norm x^2 - 2\operatorname{Re}\langle
Ux, x\rangle$ dan keuniterannya)}, lalu turunkan
$\overline{\operatorname{im}(U - I)} = F^\perp$.
(b) Tunjukkan bahwa $A_nx \to x$ untuk $x \in F$, dan $A_nx \to 0$
untuk $x \in \operatorname{im}(U - I)$ \emph{(lewat peneleskopan)},
lalu untuk $x \in \overline{\operatorname{im}(U - I)}$
(dengan batas seragam $\vertiii{A_n} \leq 1$).
(c) Simpulkan: $A_nx \to Px$ untuk \emph{setiap} $x \in H$ ---
sehingga rata-rata waktunya konvergen ke proyeksi pada invariannya.
(d) Uraikan hal itu untuk $H = L^2(\R/\Z)$ dan $Uf = f(\cdot +
\alpha)$ dengan $\alpha$ irasional: kenalilah $F$ (pakailah
deret Fourier, \cref{thm:b3:hilbert:fourier}) lalu turunkan
bahwa $\frac1n\sum_{k<n}f(x + k\alpha) \to \int_0^1f$ di
$L^2$: yakni pemerataan-$L^2$ rotasi irasionalnya.
\end{exercise}

\section{Soal: polinomial ortogonal}

\begin{problem}[{Soal akhir pekan --- Legendre, Hermite, dan
kuadratur Gauss}]\label{pb:b3:hilbert:1}
Misalkan $I \subseteq \R$ sebuah selang dan $w > 0$ sebuah
\emph{bobot} \omterm{def:b3:topology:continuity}{kontinu} pada \omterm{def:b3:topology:interior}{interior} $I$ sedemikian sehingga $\int_I
\abs t^nw(t)\dd t < \infty$ untuk setiap $n$; lalu bekerjalah di $H = L^2(I,
w\,\dd\lambda)$ dengan $\langle f, g\rangle = \int_I \bar
fg\,w$. Gram--Schmidt yang diterapkan pada $1, t, t^2, \dots$ menghasilkan
\emph{\omterm{pb:b3:hilbert:1}{polinomial ortogonal}} $(p_n)$ untuk $w$ (dengan penormalan
monik: $p_n = t^n + \cdots$).\index{polinomial
ortogonal}

\medskip
\noindent\textbf{Bagian I --- Teori umumnya.}
\begin{enumerate}
  \item Tunjukkan bahwa $p_n$ ortogonal terhadap setiap polinomial berderajat
        $< n$, dan bahwa $(p_0, \dots, p_n)$ merupakan basis
        $\R_n[t]$.
  \item (Rekurensi tiga suku) Tunjukkan bahwa ada real $a_n,
        b_n$ dengan
        \[
        p_{n+1}(t) = (t - a_n)\,p_n(t) - b_n\,p_{n-1}(t),
        \qquad b_n = \frac{\norm{p_n}^2}{\norm{p_{n-1}}^2} >
        0 .
        \]
        \emph{(Uraikan $t\,p_n$ pada basis $(p_k)_{k \leq
        n+1}$ lalu bunuhlah koefisiennya lewat keortogonalannya, dengan memakai
        $\langle tp_n, p_k\rangle = \langle p_n,
        tp_k\rangle$.)}
  \item (Akar) Tunjukkan bahwa $p_n$ mempunyai $n$ akar yang \emph{berbeda},
        semuanya \omterm{def:b3:topology:interior}{interior} terhadap $I$. \emph{(Misalkan $t_1 < \dots < t_m$
        pergantian tanda $p_n$ di \omterm{def:b3:topology:interior}{interiornya}; lalu jika $m < n$,
        ujilah $p_n$ terhadap $\prod_{i\leq m}(t - t_i)$ lalu
        pertentangkan dengan keortogonalannya.)}
\end{enumerate}

\medskip
\noindent\textbf{Bagian II --- Legendre ($I = \intcc{-1}1$, $w =
1$).} Definisikan $P_n(t) = \frac{1}{2^nn!}\,\frac{\dd^n}{\dd
t^n}\bigl[(t^2 - 1)^n\bigr]$ (Rodrigues).
\begin{enumerate}[resume]
  \item Tunjukkan bahwa $\deg P_n = n$ dengan koefisien utama
        $\frac{(2n)!}{2^n(n!)^2}$, dan, dengan mengintegralkan secara parsial
        $n$ kali, bahwa $\langle P_n, Q\rangle = 0$ untuk setiap
        polinomial $Q$ berderajat $< n$: sehingga $P_n$ merupakan (sampai
        penormalannya) \omterm{pb:b3:hilbert:1}{polinomial ortogonal} untuk $w = 1$.
  \item Hitunglah $\norm{P_n}_2^2 = \frac{2}{2n+1}$
        \emph{(integralkan secara parsial $n$ kali terhadap dirinya lalu
        susutkan ke sebuah integral Beta/Wallis,
        \cref{exo:b3:product:8})}.
  \item Tunjukkan bahwa polinomial Legendre yang ternormalkan membentuk
        \omterm{def:b3:hilbert:onb}{basis Hilbert} $L^2(\intcc{-1}1)$
        \emph{(lewat Weierstrass, \cref{cor:b3:complete:weierstrass},
        ditambah \omterm{ex:b3:lebesgue:gamma}{kepadatan} $\mathcal C$ di $L^2$)}, lalu uraikan
        $f(t) = \abs t$ sampai derajat $2$: yakni hitunglah penghampiran
        kuadratik-$L^2$ terbaik atas $\abs t$.
\end{enumerate}

\medskip
\noindent\textbf{Bagian III --- Hermite ($I = \R$, $w(t) =
\eu^{-t^2}$).} Definisikan $H_n(t) =
(-1)^n\eu^{t^2}\frac{\dd^n}{\dd t^n}\eu^{-t^2}$.
\begin{enumerate}[resume]
  \item Tunjukkan bahwa $H_n$ merupakan polinomial berderajat $n$ dengan
        koefisien utama $2^n$, bahwa $H_{n+1} = 2tH_n -
        H_n'$, dan bahwa $\langle H_m, H_n\rangle_w =
        \delta_{mn}\,2^nn!\sqrt\pi$ \emph{(lewat pengintegralan parsial lagi)}.
  \item Tunjukkan bahwa keluarga Hermitenya total di $L^2(\R,
        \eu^{-t^2}\dd t)$, dengan menerima satu hasil dari
        \cref{ch:b3:fouriertransform}: bahwa jika $g \in L^1(\R)$ mempunyai
        $\int g(t)\eu^{-\iu\xi t}\dd t = 0$ untuk setiap $\xi$,
        maka $g = 0$ hampir di mana-mana. \emph{(Untuk $f \perp$ semua $H_n$,
        yakni $\perp$ semua polinomial: tunjukkan bahwa $z \mapsto \int
        f(t)\eu^{-t^2}\eu^{-\iu zt}\dd t$ terdefinisi dengan baik,
        uraikan eksponensialnya menjadi deret, benarkan
        penukarannya lewat dominasi, lalu simpulkan bahwa
        transformasi Fourier $f\eu^{-t^2}$ lenyap.)}
\end{enumerate}

\medskip
\noindent\textbf{Bagian IV --- Kuadratur Gauss.} Tetapkan $n$, misalkan
$t_1 < \dots < t_n$ akar $p_n$ (Bagian I), lalu definisikan
bobotnya $w_i = \int_I \ell_i(t)\,w(t)\dd t$ dengan
$\ell_i$ polinomial basis interpolasi Lagrange di
$t_i$.
\begin{enumerate}[resume]
  \item Tunjukkan bahwa aturan kuadratur $Q(f) = \sum_iw_if(t_i)$
        bersifat persis pada setiap polinomial berderajat $\leq n - 1$
        (lewat interpolasinya), dan bahkan --- inilah keajaibannya --- pada
        setiap polinomial berderajat $\leq 2n - 1$: tulislah $P =
        qp_n + r$ lalu pakailah keortogonalannya pada hasil bagi $q$.
  \item Tunjukkan bahwa bobotnya positif
        \emph{(terapkan aturannya pada $\ell_i^2$, yang berderajat $2n -
        2$)}, lalu turunkan dari teorema Polya
        (\cref{exo:b3:banach:9}) bahwa kuadratur Gauss
        konvergen: yakni $Q_n(f) \to \int_I fw$ untuk setiap $f$ yang
        \omterm{def:b3:topology:continuity}{kontinu} pada $I$ yang \omterm{def:b3:topology:compact}{kompak}.
  \item Untuk $n = 2$, $I = \intcc{-1}1$, dan $w = 1$: hitunglah
        simpulnya $\pm\frac1{\sqrt3}$ dan bobotnya $1, 1$, lalu
        periksalah kepersisannya pada $1, t, t^2, t^3$ dengan tangan. Bandingkan
        dengan aturan trapesium pada kedua titik penilaian yang sama.
\end{enumerate}

\medskip
\noindent\textbf{Bagian V --- Chebyshev: polinomial yang
berosilasi paling baik.} Kini $I = \intcc{-1}1$ dan $w(t) =
\frac1{\sqrt{1 - t^2}}$.
\begin{enumerate}[resume]
  \item Tunjukkan bahwa $T_n(\cos\theta) = \cos n\theta$ mendefinisikan sebuah
        polinomial $T_n$ berderajat $n$ (tegakkanlah $T_{n+1} =
        2t\,T_n - T_{n-1}$ dari sebuah identitas trigonometri),
        dengan koefisien utama $2^{n-1}$ untuk $n \geq 1$;
        dan bahwa substitusi $t = \cos\theta$ memberikan
        \[
        \langle T_m, T_n\rangle_w =
        \int_0^\pi\cos m\theta\,\cos n\theta\,\dd\theta
        = 0 \ (m \neq n), \qquad
        \norm{T_0}_w^2 = \pi,\ \ \norm{T_n}_w^2 = \frac\pi2 :
        \]
        sehingga $T_n$ merupakan \omterm{pb:b3:hilbert:1}{polinomial ortogonal} untuk bobot
        ini, dan uraian Chebyshevnya \emph{adalah} deret kosinus
        Fourier yang menyamar.
  \item Temukan secara eksplisit $n$ akarnya $t_k =
        \cos\frac{(2k-1)\pi}{2n}$ dan $n + 1$ ekstremumnya
        $s_j = \cos\frac{j\pi}n$ dari $T_n$ pada $\intcc{-1}1$,
        yang $T_n(s_j) = (-1)^j$-nya: sehingga grafiknya
        \emph{berosilasi setara} antara $\pm1$.
  \item (Minimaks) Tunjukkan bahwa di antara semua polinomial \emph{monik}
        berderajat $n$, polinomial
        $2^{1-n}T_n$ mempunyai norma-sup terkecil pada
        $\intcc{-1}1$, yakni $2^{1-n}$ --- dan ia
        satu-satunya pemberi minimumnya. \emph{(Sebab jika sebuah $P$ yang monik mempunyai
        $\sup\abs P < 2^{1-n}$, maka selisih $2^{1-n}T_n -
        P$, yang berderajat $\leq n-1$, akan berganti tanda di
        $n+1$ titik osilasi setaranya.)}
  \item Terapannya pada interpolasi: untuk simpul $t_1 < \dots
        < t_n$ di $\intcc{-1}1$, galat interpolasi
        Lagrange atas fungsi $\mathcal C^n$ melibatkan
        $\omega(t) = \prod_i(t - t_i)$. Tunjukkan bahwa memilih
        akar Chebyshev sebagai simpulnya meminimumkan
        $\sup_{\intcc{-1}1}\abs\omega$, lalu berikan
        batas yang dihasilkannya $\norm{f -
        L_nf}_\infty \leq \frac{\norm{f^{(n)}}_\infty}
        {2^{n-1}\,n!}$ --- lalu bandingkan dengan simpul berjarak sama
        (dengan menyatakan gejala Runge sebagai kisah peringatannya).
  \item Periksalah $\abs{T_n'(\pm1)} = n^2$ \emph{(turunkan
        $T_n(\cos\theta) = \cos n\theta$ lalu ambil limitnya
        $\theta \to 0, \pi$)}: sehingga polinomial yang terbatas oleh $1$ pada
        $\intcc{-1}1$ dapat berturunan sebesar $n^2$
        di tepinya (dan ketaksamaan Markov mengatakan tak lebih besar ---
        pernyataannya saja). Di manakah pada selangnya
        batas turunannya hanya $O(n)$?
  \item (Kuadratur Chebyshev--Gauss) Tunjukkan bahwa aturan Gauss
        untuk bobot $w$ pada $n$ akar Chebyshevnya
        mempunyai bobot yang \emph{sama} $w_i = \frac\pi n$
        \emph{(lewat kepersisannya pada $T_0, \dots, T_{n-1}$ ditambah
        jumlah trigonometrinya
        $\sum_{k=1}^n\cos\bigl(j\tfrac{(2k-1)\pi}{2n}\bigr) =
        0$ untuk $1 \leq j \leq n - 1$)}: yakni yang paling seragam di antara
        semua kuadraturnya. Tuliskanlah untuk $n = 3$.
\end{enumerate}

\medskip
\noindent\textbf{Bagian VI --- Christoffel--Darboux,
keberselang-selingan, dan matriks Jacobi.} Kembali ke bobot umum;
dengan $h_k = \norm{p_k}^2$ ($p_k$ monik) dan $b_k = h_k/h_{k-1}$.
\begin{enumerate}[resume]
  \item (Norma terkecil) Tunjukkan bahwa di antara semua polinomial
        \emph{monik} berderajat $n$, $p_n$ yang ortogonal itu
        satu-satunya yang bernorma-$L^2(w)$ minimal --- lalu kenalilah
        peminimumannya sebagai sebuah proyeksi ortogonal ke
        $\R_{n-1}[t]$ (\cref{thm:b3:hilbert:projection} atau
        proyeksi berdimensi berhingga pada Tahun ke-2). Sedangkan
        sifat minimaks pertanyaan 14 adalah pernyataan yang sama
        dengan $L^\infty$ menggantikan $L^2$: yakni pahlawan yang sama, dua
        norma.
  \item (Christoffel--Darboux) Buktikan, lewat induksi pada $n$
        dengan memakai rekurensi tiga sukunya, identitas
        \[
        \sum_{k=0}^{n}\frac{p_k(x)\,p_k(y)}{h_k}
        = \frac{p_{n+1}(x)\,p_n(y) -
        p_n(x)\,p_{n+1}(y)}{h_n\,(x - y)}
        \qquad (x \neq y),
        \]
        beserta bentuk konfluennya (saat $y \to x$):
        $\sum_{k\leq n}\frac{p_k(x)^2}{h_k} =
        \frac{p_{n+1}'(x)p_n(x) - p_n'(x)p_{n+1}(x)}{h_n}$.
  \item Turunkan bahwa $p_n$ dan $p_{n+1}$ tak berakar bersama,
        dan bahwa pada setiap akar $x_0$ dari $p_{n+1}$ berlaku:
        $p_n(x_0)\,p_{n+1}'(x_0) > 0$. Lalu simpulkan
        \emph{keberselang-selingan} akarnya: bahwa di antara dua akar berurutan
        $p_{n+1}$ terletak tepat satu akar $p_n$.
  \item (Matriks Jacobi) Misalkan $J_n$ matriks tridiagonal
        setangkup berukuran $n\times n$ dengan diagonal $a_0,
        \dots, a_{n-1}$ dan entri luar diagonalnya $\sqrt{b_1},
        \dots, \sqrt{b_{n-1}}$. Tunjukkan lewat induksi bahwa
        $\det(tI_n - J_n) = p_n(t)$, sehingga akar $p_n$
        merupakan nilai eigen sebuah matriks setangkup real ---
        yang membuktikan kembali dalam satu baris bahwa akarnya real, dan (bersama
        keberselang-selingan di atas) mengikat \omterm{pb:b3:hilbert:1}{polinomial ortogonalnya}
        ke dunia spektral \cref{ch:b3:spectral}.
  \item (Sintesis) Rakitlah kamus bagi ketiga
        keluarga klasiknya (Legendre, Hermite, Chebyshev):
        yakni selangnya, bobotnya, rumus pendefinisinya, rekurensi
        tiga sukunya, normanya, dan habitat alami masing-masingnya
        (kuadratur dan penghampiran pada \omterm{def:b3:topology:compact}{kompak}; analisis Gauss;
        serta metode minimaks dan kosinus Fourier). Lalu satu
        kalimat tentang apa yang diberikan teori umumnya (Bagian I, VI)
        yang tak dapat diberikan perhitungan satuan mana pun.
\end{enumerate}

\medskip
\noindent\textbf{Bagian VII --- Suku galatnya, dan kernel
di balik bobotnya.} Di sini $I$ \omterm{def:b3:topology:compact}{kompak} dan $f \in \mathcal
C^{2n}(I)$.
\begin{enumerate}[resume]
  \item (Rumus galat Gauss) Misalkan $Hf$ interpolan Hermite
        berderajat $\leq 2n - 1$ yang mencocokkan $f$ dan
        $f'$ di simpul $t_1, \dots, t_n$ (buktikanlah
        keberadaannya dan galat titik demi titiknya
        \[
        f(t) - Hf(t) =
        \frac{f^{(2n)}(\xi_t)}{(2n)!}\;p_n(t)^2
        \]
        lewat argumen fungsi bantunya yang lazim). Lalu turunkan, dengan
        mengintegralkan identitas ini terhadap $w$ dan menjepitnya
        di antara ekstremum $f^{(2n)}$, bahwa
        \[
        \int_I f\,w - Q_n(f)
        = \frac{f^{(2n)}(\xi)}{(2n)!}\,h_n
        \qquad\text{untuk suatu } \xi \in I,
        \]
        dengan $h_n = \norm{p_n}^2$ seperti pada Bagian VI: sehingga kuadratur
        Gauss meleset sebesar satu turunan ke-$2n$, yang berbobot
        norma kuadrat \omterm{pb:b3:hilbert:1}{polinomial ortogonal} moniknya.
  \item (Bobotnya adalah nilai Christoffel) Dengan memakai
        kernel pereproduksi $K_n(x, y) = \sum_{k=0}^{n-1}
        \frac{p_k(x)p_k(y)}{h_k}$ dari $\R_{n-1}[t]$ dan
        kepersisan $Q_n$ sampai derajat $2n - 2$, buktikan bahwa
        \[
        w_i \;=\;
        \Bigl(\,\sum_{k=0}^{n-1}
        \frac{p_k(t_i)^2}{h_k}\Bigr)^{\!-1} :
        \]
        yakni setiap bobotnya merupakan nilai di simpulnya bagi
        \emph{fungsi Christoffel} --- jadi kepositifan
        bobotnya (pertanyaan 10) lagi, kini dengan rumus yang
        persis. Periksalah bahwa ia memulihkan $w_1 = w_2 = 1$ untuk $n =
        2$, $I = \intcc{-1}1$, dan $w = 1$.
  \item (Segalanya cocok pada satu integral) Untuk bobot Chebyshev
        dan $n = 3$ simpul, hitunglah kedua ruas
        \[
        \int_{-1}^{1}\frac{t^6}{\sqrt{1 - t^2}}\,\dd t
        = \frac{5\pi}{16},
        \qquad
        Q_3(t^6) = \frac{9\pi}{32},
        \]
        sehingga galat kuadraturnya tepat $\frac{\pi}{32}$;
        lalu periksalah bahwa rumus galat pertanyaan 23
        meramalkan persis nilai itu (sebab di sini $f^{(6)} = 6!$
        konstan, dan $h_3 = \norm{2^{-2}T_3}_w^2 =
        \frac\pi{32}$): sehingga teori dan perhitungannya bersesuaian sampai
        angka terakhirnya.

\end{enumerate}
\end{problem}
